\documentclass[11pt,a4paper]{article}  % {{{*
\usepackage[utf8]{inputenc}
\usepackage[spanish,es-tabla]{babel} % 'es-tabla' cambia "Cuadro" por "Tabla"
\usepackage{amsmath}
\usepackage{amssymb}
\usepackage{geometry}
\usepackage[normalem]{ulem}
\usepackage{xparse} % Necesario para argumentos opcionales delimitados con { } (tipo g)

\geometry{margin=2.5cm}

% --- Configuración de la bibliografía ---
\usepackage{cite} 
\usepackage{xcolor}
\usepackage[colorlinks=true, linkcolor=blue, citecolor=blue, urlcolor=blue]{hyperref}
\usepackage[draft,inline,nomargin]{fixme} \fxsetup{theme=color}
\FXRegisterAuthor{cp}{acp}{\color{blue}CP}
\FXRegisterAuthor{sn}{asn}{\color{red}SN}

\definecolor{jacolor}{HTML}{C82800}
\FXRegisterAuthor{ja}{aja}{\color{jacolor}JA}

% Notes with one argument keep the usual FiXme behavior.  With two
% arguments they show Insert/Reason; with three, Remove/Insert/Reason.
\NewDocumentCommand{\flexfxnote}{m m m m}{%
  \IfNoValueTF{#3}{#1{#2}}{%
    \IfNoValueTF{#4}{%
      #1{\sout{#2} #3}%
    }{%
      #1{\sout{#2} #3 {\small \tt #4}}%
    }%
  }%
}
\NewCommandCopy{\fxnoteoriginal}{\fxnote}
\NewCommandCopy{\janoteoriginal}{\janote}
\NewCommandCopy{\cpnoteoriginal}{\cpnote}
\NewCommandCopy{\mpnoteoriginal}{\mpnote}
\NewCommandCopy{\jnnoteoriginal}{\jnnote}
\NewCommandCopy{\clnoteoriginal}{\clnote}
\NewCommandCopy{\codexnoteoriginal}{\codexnote}
\RenewDocumentCommand{\fxnote}{m g g}{%
  \flexfxnote{\fxnoteoriginal}{#1}{#2}{#3}}
\RenewDocumentCommand{\janote}{m g g}{%
  \flexfxnote{\janoteoriginal}{#1}{#2}{#3}}
\RenewDocumentCommand{\cpnote}{m g g}{%
  \flexfxnote{\cpnoteoriginal}{#1}{#2}{#3}}

\newcommand{\jainsert}[1]{\janote{#1}}
\newcommand{\jaremove}[1]{\janote{\sout{#1}}}
\newcommand{\jareason}[1]{\janote{\tt #1}}
\newcommand{\jacomment}[2]{\uline{#1}\janote{#2}}
  % *}}}
\begin{document}
% Título y otras cosas {{{*
\begin{center}
    \Large \textbf{UNIVERSIDAD DE SAN CARLOS DE GUATEMALA} \\
    \textbf{ESCUELA DE CIENCIAS FÍSICAS Y MATEMÁTICAS} \\[1.5cm]
    \Large \textbf{Protocolo de trabajo de tesis:} \\
    \textbf{Caminatas cuánticas con una moneda en 4D} \\[1.5cm]
    \textbf{Saúl Estuardo Nájera Allara} \\
    Carné: 202107506 \\[1cm]
    Septiembre 2026 \\[1cm]
    Licenciatura en Física Aplicada \\[1.5cm]
    \textbf{Asesorado por:} \\
    M.Sc. José Alfredo de León (IF-UNAM), \\
    Dr. Carlos Francisco Pineda Zorrilla (IF-UNAM), \\
    Ing. Rodolfo Samayoa (ECFM-USAC)
\end{center}
  % *}}}
\newpage \tableofcontents \newpage
\section{Introducción}  % {{{*

En la literatura reciente, se ha estudiado un modelo de caminata cuántica bidimensional que emplea un solo qubit como estado de moneda~\cite{alonso-lobo2025simplest}, 
concluyendo que existe una correspondencia entre las condiciones de frontera y la dinámica del sistema.
Sin embargo, no es concluyente si esta relación se mantiene al emplear un estado de moneda conformado por un único sistema de cuatro dimensiones (qudit).
Clásicamente, una partícula libre en un dominio bidimensional presenta una dinámica caótica o integrable que depende de la geometría de sus fronteras. Motivado por esto, el presente trabajo propone investigar si la dinámica de un caminante con una moneda en 4D exhibe un comportamiento análogo. Específicamente, se buscará determinar si las condiciones de frontera por sí solas son suficientes para establecer la caoticidad o integrabilidad del sistema.

Este proyecto se desarrolla en el marco del Instituto de investigación de Ciencias Físicas y Matemáticas (ICFM) de la USAC, en colaboración con el Instituto de Física de la UNAM. El trabajo se lleva a cabo en conjunto con el grupo de Información y Óptica cuántica de la UNAM (GIOC-UNAM). La investigación está bajo la asesoría del M.Sc. José Alfredo de León, fungiendo como cotutores el Dr. Carlos Francisco Pineda Zorrilla y el Ing. Rodolfo Samayoa.
% *}}}
\section{Objetivos}  % {{{*

\subsection{Objetivo general}
Caracterizar la emergencia de caos cuántico en un modelo de caminata cuántica discreta en el tiempo con una moneda de cuatro dimensiones y un espacio de posición con geometría de billar. 

\subsection{Objetivos específicos}
\begin{itemize}
    \item Implementar numéricamente el operador de evolución de un paso para el modelo de caminata cuántica discreta en el tiempo en espacios de posición acotados por billares con geometrías de rectángulo y Sinaí.
    \item Determinar la densidad de cuasienergías del operador de evolución para identificar singularidades espectrales o regularidades estructurales asociadas a cada geometría.
    \item Contrastar las predicciones de la teoría de matrices aleatorias con los resultados numéricos de la distribución del espaciamiento de niveles a primeros vecinos y el factor de forma espectral para determinar la existencia de correlaciones de corto y largo alcance en el espectro de cuasienergías.
    \item Calcular la evolución temporal de la razón inversa de participación (IPR) en el espacio de posición para buscar las señales de (des)localización típicas de sistemas caóticos. 
\end{itemize}
% *}}}
\section{Marco teórico}  % {{{* 

Las caminatas cuánticas son el análogo cuántico de una caminata aleatoria clásica. El estado de un caminante es el producto tensorial de su posición y el estado de su moneda.
En este trabajo se considerará un dominio espacial bidimensional y un espacio de moneda de dimensión 4. 
La evolución del caminante está dictada por la aplicación del operador unitario de evolución $\mathcal{U}$. Basándose en la literatura existente~\cite{portugal2018quantum}, una forma posible de definir dicho operador es:
\begin{equation}
    \mathcal{U}=S\cdot C.
    \label{eq:U}
\end{equation}
Siendo $S$ un operador que desplaza al caminante según el estado de su moneda y $C$ un operador que actúa únicamente sobre el estado de la moneda.
La forma de $S$ depende de las condiciones de frontera a considerar y la evolución de un estado según \eqref{eq:U} está influenciada por la elección del operador $C$.
Las caminatas cuánticas proporcionan un marco teórico para el diseño de nuevos algoritmos cuánticos~\cite{shenvi2003quantum},
por ejemplo, el algoritmo de Grover~\cite{grover1996fastquantummechanicalalgorithm} para la búsqueda en bases de datos no estructuradas puede reformularse como una caminata cuántica sobre un grafo completo~\cite{portugal2018quantum}.

El estudio del caos cuántico se originó a partir del análisis de sistemas cuánticos cuyo límite clásico exhibe una dinámica caótica. Esta investigación condujo a establecer una correspondencia entre la estadística de los espectros de energía de dichos sistemas y la estadística de los ensambles descritos por la teoría de matrices aleatorias (RMT, por sus siglas en inglés)~\cite{bohigas1984characterization}.
En este marco, la conjetura de Berry-Tabor postula que las fluctuaciones espectrales de sistemas con una contraparte clásica integrable obedecen una estadística de Poisson~\cite{berry1977level}. Por su parte, la conjetura de Bohigas, Giannoni y Schmit (BGS) establece que, si el análogo clásico es caótico, las propiedades estadísticas del espectro seguirán de manera universal las predicciones de RMT~\cite{bohigas1984characterization}.
Para aplicar este formalismo a las caminatas cuánticas, se debe considerar que el operador de evolución $\mathcal{U}$ es unitario y sus valores propios son complejos de norma 1; por lo tanto, la presencia de caos cuántico no se evalúa con los valores propios de \eqref{eq:U}, sino sobre las fases de dichos valores propios, conocidas como cuasienergías, analizando sus correlaciones estadísticas tanto de corto como de largo alcance.

Para estudiar las correlaciones de corto alcance se emplean las diferencias a primeros vecinos, $s$, definida como la separación entre dos cuasienergías consecutivas normalizada por el espaciamiento medio del espectro. 
Haciendo uso de las predicciones analíticas de RMT~\cite{wigner1951statistical,mehta1991random}, se observa que la densidad de probabilidad $P(s)$ para sistemas caóticos exhibe repulsión de niveles ($P(s)\rightarrow 0,\, s\rightarrow 0$). En contraste, los sistemas integrables obedecen una estadística de Poisson ($P(s) = e^{-s}$), esta distribución produce una máxima probabilidad para espaciamientos pequeños, fenómeno que se visualiza como una aglomeración de niveles o \textit{level clustering}~\cite{mehta1991random,alonso-lobo2025simplest}.

Por otro lado, para detectar correlaciones de largo alcance se empleará el factor de forma espectral (SFF por sus siglas en inglés)~\cite{haake2018quantum}. La manifestación de caos cuántico para este indicador consiste en la aparición de un agujero de correlaciones (\textit{dip}), seguido de una rampa (\textit{ramp}) y una saturación alrededor de un valor constante (\textit{plateau}), características que se revelan únicamente después de aplicar un promedio estadístico adecuado (sobre ventanas móviles o ensambles)~\cite{haake2018quantum}.
Si el sistema es integrable, no existe un \textit{dip} ni una \textit{ramp}, presentándose únicamente el \textit{plateau}. 
% *}}}
\section{Metodología}  % {{{*
\snnote{Siendo revisada por: CP}
\cpnote{Listo}

En primer lugar, la implementación del operador de evolución unitaria dado por la ecuación~\eqref{eq:U} y el análisis del sistema se llevarán a cabo empleando el paquete especializado \textit{QuantumWalks}~\cite{quantumwalks} desarrollado en Wolfram Mathematica, del cual el autor del presente trabajo es contribuyente. 
Para ello, haciendo uso de este paquete se construirán los espacios de posición acotados por billares con geometrías de rectángulo y Sinaí. La estructura del operador de desplazamiento $S$ se generará mediante reglas que condicionan el movimiento del caminante hacia sus primeros vecinos siguiendo un esquema de cruz (es decir, permitiendo únicamente transiciones a lo largo de los ejes cartesianos de la red bidimensional), según el estado de la moneda.
Por su parte, la dinámica interna del caminante estará descrita mediante el operador $C = \mathbb{I}_P \otimes \mathcal{C}$, donde $\mathbb{I}_P$ representa el operador identidad en el espacio de posición y $\mathcal{C}$ es un operador unitario que actúa exclusivamente sobre el espacio de la moneda de cuatro dimensiones.


A partir del operador~\eqref{eq:U}, el siguiente paso consiste en extraer su
espectro de valores propios. Dado que la evolución es descrita por un operador
unitario, dichos valores son números complejos de norma 1. El análisis se
enfocará exclusivamente en las fases de estos valores, las cuales corresponden
a las cuasienergías del sistema y servirán como base para el estudio
estadístico. Sin embargo, previo al estudio estadístico, se aplicará un
procedimiento de desdoblamiento espectral (\textit{unfolding}) a este conjunto
de cuasienergías para separar la densidad media de estados de las fluctuaciones
espectrales locales, garantizando así un espaciamiento medio unitario que
permita la comparación directa con las predicciones de RMT.

Con el espectro de cuasienergías obtenido, se procederá a evaluar las
manifestaciones de caos cuántico contrastando los resultados numéricos con las
predicciones de RMT. Para evaluar las correlaciones de corto alcance, se
calculará la distribución del espaciamiento a primeros vecinos $P(s)$, buscando
identificar repulsión de niveles o estadística de Poisson. Adicionalmente, para
detectar correlaciones de largo alcance, se calculará el factor de forma
espectral (SFF) aplicando los promedios estadísticos apropiados para visualizar
la posible formación de un agujero de correlaciones, una rampa y su saturación.

Para complementar el análisis espectral, se estudiarán las consecuencias
espaciales de la dinámica. En primer lugar, se implementará una herramienta de
visualización de la función de onda en el espacio de posición; esto se logrará
calculando la probabilidad de hallar al caminante en cada nodo de la red para
un tiempo $t$, independientemente de su estado de moneda. A partir de esta
distribución de probabilidad espacial, se procederá a calcular la evolución
temporal de la razón inversa de participación (IPR). El comportamiento conjunto
de estas herramientas permitirá visualizar y caracterizar formalmente las
señales de localización o deslocalización espacial, fenómenos que diferencian
típicamente a los regímenes caóticos de los integrables.

% *}}}
\section{Justificación}  % {{{*

\snnote{Siendo revisada por: JA}\janote{Tiene mi VoBo}

Las caminatas cuánticas poseen una relevancia científica importante. Por un lado, son la base para la implementación de nuevos algoritmos cuánticos~\cite{shenvi2003quantum}. Por otro lado, al constituir un modelo de computación universal~\cite{childs2009universal,lovett2010universal}, permiten simular una amplia variedad de sistemas físicos complejos.
Esto facilita el estudio de fenómenos como el transporte coherente\cpnote{Para cada uno de los fenomenos, tienes que dar una o dos citas}\snnote{He añadido referencias representativas de la literatura para cada fenómeno.}~\cite{mulken2011continuous}, las transiciones de fase topológicas~\cite{kitagawa2010exploring} y las dinámicas en redes estructuradas~\cite{kempe2003quantum,portugal2018quantum}.
Además, estos modelos pueden ser implementados directamente en diversas plataformas experimentales, como circuitos fotónicos\cpnote{cita}~\cite{peruzzo2010quantum} o redes ópticas\cpnote{cita, por favor cuando pongas las citas, trata de que esté el link y que se pueda hacer click, para verificar que si citaste bien}\snnote{Añadidas citas.}~\cite{karski2009quantum}. 
Esta viabilidad experimental representa una ventaja, ya que permite investigar fenómenos cuánticos evadiendo las limitaciones tecnológicas, como la decohrencia, que actualmente impiden la construcción de computadoras cuánticas universales a gran escala~\cite{portugal2018quantum,lorz2019photonic}. 
Gracias a estos montajes experimentales, los investigadores disponen de un entorno controlable~\cite{kendon2003decoherence} para explorar la dinámica de la información de los sistemas a trabajar.

Uno de los fenómenos que pueden ser explorados aprovechando estos entornos controlables es el caos cuántico. 
Para estudiar este fenómeno de forma aislada, los sistemas de billar resultan ideales, ya que clásicamente permiten transicionar de una dinámica integrable a una caótica alterando únicamente la geometría de sus fronteras, como ocurre al pasar de un rectángulo a un billar de Sinaí~\cite{sinai1970dynamical}.
Para investigar cómo se manifiestan estas propiedades en el régimen cuántico, 
las caminatas cuánticas discretas en el tiempo (DTQWs por sus siglas en inglés) 
ofrecen un marco de estudio ventajoso, ya que nos permiten variar de manera 
independiente tanto las condiciones de frontera como el operador de moneda, 
ofreciendo así las herramientas necesarias para determinar si la geometría de las fronteras es, por sí sola, suficiente para dar lugar a la emergencia del caos.

Aunque este enfoque geométrico ha comenzado a explorarse, evidenciando una correspondencia preliminar entre las condiciones de frontera y la dinámica espectral en DTQWs bidimensionales,
dichos trabajos se han restringido al uso de un solo sistema de dos niveles como subsistema de moneda~\cite{alonso-lobo2025simplest}. Esta limitación presenta dos problemas importantes. 
Por un lado, las distribuciones espectrales obtenidas bajo este enfoque no concuerdan exactamente con las predicciones universales del Ensamble Ortogonal Gaussiano (GOE) esperadas para sistemas caóticos. 
Por otro lado, dado que un caminante en una red bidimensional cuenta con cuatro direcciones posibles de desplazamiento (arriba, abajo, derecha, izquierda), 
emplear un sistema de dos niveles (qubit) para gobernar esta dinámica resulta estructuralmente restrictivo y artificial. Resulta mucho más natural y directo codificar estos cuatro movimientos ortogonales en un sistema de moneda que posea la misma dimensionalidad, es decir, un estado de cuatro dimensiones (qudit). Por lo tanto, es de interés investigar si esta formulación permite que el comportamiento espectral siga las predicciones estadísticas esperadas, permitiendo determinar si la geometría de las fronteras es una condición suficiente para dictar, por sí sola, el caos o la integrabilidad del sistema.

% *}}}
% --- Bibliografía ---  % {{{*
\bibliographystyle{unsrt} % Estilo de la bibliografía (orden de aparición)
\bibliography{ref}        % Nombre de tu archivo .bib (sin la extensión)
% *}}}

\end{document}
